\section{Background}

\subsection{The Capacity-Conditional Calibration Problem}

Papers~7~\cite{paper7online} and~8~\cite{paper8multi} together
established that no fixed-history calibration strategy is universally
deployable for HygienePrio's scorer weights across capacity regimes:
lag-1 grid search recovers the offline-peek ceiling at $K \in
\{50, 100\}$ but harms performance at $K = 200$, and multi-window
smoothing (EWMA-3, trailing-mean-3) amplifies the harm rather than
reversing it. Paper~8 explicitly named adaptive procedures ---
change-point detection or capacity-conditional gating --- as the
remaining candidate class.

\subsection{Change-Point Detection}

A change-point detector flags windows at which the underlying
distribution has shifted enough that recently-fit parameters are
unlikely to generalise to the current observation. Classical
examples include CUSUM~\cite{page1954cusum} and Bayesian online
change-point methods; both are out of scope for this paper. The
detector studied here is intentionally minimal: a single-threshold
magnitude comparison on the calibration-target P@50 stream.

\subsection{Capacity-Conditional Gating}

A static alternative to adaptive switching is a capacity-conditional
gate: deploy lag-1 below a threshold capacity and fixed weights above
it. This requires no per-window decision but assumes a clean
capacity-boundary, which Paper~7's results gesture toward (clean
recovery at $K \leq 100$, clear harm at $K = 200$). Whether an
adaptive detector outperforms this static rule is RQ4.

\subsection{Why a Minimal Detector First}

Before evaluating sophisticated detectors, it is useful to know whether
even the simplest plausible detector beats both lag-1 and fixed in
the regimes where each individually fails. If a one-threshold rule
suffices, more elaborate machinery is not needed; if it fails, the
failure pattern indicates where the elaborate machinery is required.
